% Continuacion de la introduccion (\subsection). Contexto geografico de las
% estaciones: la rubrica pide describir ubicacion y factores geograficos que
% podrian influir en las mediciones. Coordenadas y elevacion vienen del
% catalogo docs/Ubicacion de las estaciones de monitoreo.docx; la lectura
% espacial, de docs/05_estacionalidad_ozono.pdf.

\subsection{Las estaciones y su contexto geográfico}
\label{sec:contexto-estaciones}

El AMM ocupa un valle a unos 500 msnm limitado al sur y al poniente por la
Sierra Madre Oriental, al oriente por el Cerro de la Silla y al norte por el
Cerro del Topo Chico y la Sierra de las Mitras. Las 15 estaciones del SIMA se
reparten en ese valle a lo largo de un gradiente de elevación de 334 msnm,
en la planicie del sureste, a 702 msnm, al pie de la sierra en el noroeste
(\cref{tab:estaciones}). El código de cada estación indica su cuadrante, y a
partir de él se homologan en siete zonas que los modelos de
\cref{sec:clasificacion_tecs} usan como predictor categórico.

\begin{table}[t]
  \centering
  \caption{Estaciones del SIMA, zona a la que se homologan y elevación.
    Coordenadas en grados y minutos; todas en latitud norte y longitud
    oeste.}
  \label{tab:estaciones}
  \small
  \begin{tabular}{llrll}
    \toprule
    Código & Zona & msnm & Lat. & Long. \\
    \midrule
    \texttt{SE3}  & Sureste  & 334 & 25°36' & 100°00'  \\
    \texttt{NE3}  & Noreste  & 346 & 25°47' & 100°05' \\
    \texttt{SE2}  & Sureste  & 387 & 25°39' & 100°06' \\
    \texttt{NE2}  & Noreste  & 432 & 25°47' & 100°11' \\
    \texttt{NE}   & Noreste  & 474 & 25°45' & 100°15' \\
    \texttt{SE}   & Sureste  & 500 & 25°40' & 100°15' \\
    \texttt{NTE}  & Norte    & 503 & 25°48' & 100°20' \\
    \texttt{NTE2} & Norte    & 520 & 25°44' & 100°19' \\
    \texttt{SUR}  & Sur      & 555 & 25°37' & 100°16' \\
    \texttt{CE}   & Centro   & 562 & 25°41' & 100°20' \\
    \texttt{NO}   & Noroeste & 568 & 25°46' & 100°22' \\
    \texttt{NO3}  & Noroeste & 607 & 25°46' & 100°28' \\
    \texttt{SO2}  & Suroeste & 636 & 25°40' & 100°25' \\
    \texttt{SO}   & Suroeste & 674 & 25°41' & 100°27' \\
    \texttt{NO2}  & Noroeste & 702 & 25°48' & 100°35' \\
    \bottomrule
  \end{tabular}
\end{table}

Tres factores geográficos condicionan las mediciones. El primero es el
viento: la mediana de la componente zonal del viento es negativa (\cref{tab:cuartiles-iqr}), es decir, el
viento predominante entra del oriente y arrastra los precursores emitidos en
el centro y el oriente de la ciudad hacia el poniente, donde la sierra frena
su salida. El segundo es la elevación, que covaría con la distancia a las
fuentes de \ce{NOx} y con el uso de suelo: la correlación entre la elevación
de una estación y su tasa de excedencia es 0.827 sobre las 15 estaciones
(\cref{fig:elevacion-excedencia}). El tercero es el encierro orográfico ya
descrito en \cref{sec:objetivo}, que en episodios de alta presión y viento
débil impide la dispersión vertical y horizontal. El resultado es un
ordenamiento poniente--oriente estable en las cuatro temporadas: el Noroeste
y el Suroeste concentran las tasas de excedencia más altas (43.3\,\% y
42.1\,\%) y el Noreste la más baja (19.5\,\%).

\begin{figure}[t]
  \centering
  \includegraphics[width=\linewidth]{elevacion_vs_excedencia.pdf}
  \caption{Tasa de excedencia del MDA8 a 51 ppb contra elevación de cada
    estación, 2021--2025. Con $n=15$ y estaciones vecinas correlacionadas, la
    recta describe una asociación, no un efecto de la altitud.}
  \label{fig:elevacion-excedencia}
\end{figure}

Dos consideraciones acotan la comparabilidad entre estaciones.
\texttt{NO3} opera desde diciembre de 2022 y \texttt{NE3} entró en operación
después que el resto, de modo que su ventana es más corta; y ninguna estación
mide compuestos orgánicos volátiles, el segundo precursor del ozono, lo que
limita toda interpretación química a los óxidos de nitrógeno y al \ce{CO}
como sustituto.
